%!TEX root = ../../document/document.tex

% Only used acronyms will be displayed in the list of acronyms.
% For more options see: http://www.namsu.de/Extra/pakete/Acronym.pdf

% Usage:
% \acro{BSP}{Board Support Package} % Example acronym
% \ac{Abbr.}   --> inserts the abbreviation, at the first call the full version is inserted
% \acf{Abbr.}   --> inserts the abbreviation AND the explanation
% \acl{Abbr.}   --> inserts only the explanation
% \acp{Abbr.}  --> inserts the plural form (adds 's'); the additional 'p' works with the above commands as well
% \acs{Abbr.}   --> inserts the abbreviation

% If the space between the acronym and the explanation is too small, you can adjust it by going into the config/additionals/acronyms.tex file and changing the 'ABCDEFGHIJ' to the longest acronym you have. 


%A:
\acro{AI}{Artificial Intelligence}
\acro{API}{Application Programming Interface}

%B:

%C:
\acro{CI}{Continuous Integration}
\acro{CD}{Continuous Deployment}

%D:
\acro{dta}{das Test-Acronym}
\acro{DHBW}{Duale Hochschule Baden-Württemberg}

%E:

%F:

%G:

%H:
\acro{HTML}{Hypertext Markup Language}

%I:
\acro{IaaS}{Infrastructure as a Service}
\acro{IaC}{Infrastructure as Code}
\acro{IP}{Internet Protocol}
\acro{IPA}{\acl{IP} Address}
\acro{IT}{Information Technology}
\acro{PaaS}{Platform as a Service}
\acro{SaaS}{Software as a Service}

%J:
\acro{JSON}{JavaScript Object Notation}

%K:

%L:

%M:

%N:
\acro{NSG}{Network Security Group}

%O:

%P:
\acro{PC}{Personal Computer}

%Q:

%R:

%S:
\acro{SSH}{Secure Shell}

%T:

%U:

%V:
\acro{VM}{Virtual Machine}

%W:
\acro{WSGI}{Web Server Gateway Interface}

%X:

%Y:
\acro{YAML}{Yet Another Markup Language}

%Z:
